\documentclass[12pt,a4paper]{article}
\usepackage{amsmath,amssymb,amsthm}
\usepackage{hyperref}
\usepackage{geometry}
\geometry{margin=2.5cm}
\usepackage{enumitem}
\usepackage{booktabs}

\newtheorem{theorem}{Theorem}[section]
\newtheorem{lemma}[theorem]{Lemma}
\newtheorem{corollary}[theorem]{Corollary}
\newtheorem{proposition}[theorem]{Proposition}
\theoremstyle{definition}
\newtheorem{definition}[theorem]{Definition}
\theoremstyle{remark}
\newtheorem{remark}[theorem]{Remark}

\title{Quantum Mechanics Derived from Cost Minimization:\\
The Forcing Chain from a Single Primitive to Full Quantum Theory}

\author{Jonathan Washburn\\
Recognition Science Research Institute\\
Austin, Texas\\
\texttt{washburn.jonathan@gmail.com}}

\date{March 2026}

\begin{document}
\maketitle

\begin{abstract}
We derive the full structure of quantum mechanics from a single
primitive: the Recognition Composition Law (RCL).  A nine-level forcing
chain (T0--T8) produces, in succession: logic (T0), existence over
non-existence (T1), discreteness (T2), double-entry bookkeeping (T3),
the recognition operator $\widehat{R}$ (T4), the unique cost functional
$J(x) = \tfrac{1}{2}(x + x^{-1}) - 1$ (T5), the golden ratio
$\varphi = (1+\sqrt{5})/2$ (T6), the 8-tick epoch (T7), and spatial
dimension $D = 3$ (T8).  From these forced structures, quantum
mechanics emerges: complex Hilbert space (from the 8-tick half-period
$e^{i\pi} = -1$), the Born rule $P = |\psi|^2$ (from $J$-cost phase
invariance), entanglement (from shared ledger constraints),
commutation relations $[\hat{x}, \hat{p}] = i\hbar$ (from double-entry
non-commutativity), and definite measurement outcomes (from the
discrete ledger update at $C \geq 1$).  No measurement postulate, no
Hilbert-space axiom, no Born rule assumption, and no free parameters
are required.  \end{abstract}

%% ============================================================
\section{Introduction}
%% ============================================================

Quantum mechanics has a structure that works with extraordinary
precision, yet its foundations remain disputed~\cite{Hardy2001,
Chiribella2011}.  Five foundational questions persist:
\begin{enumerate}
  \item Why the Born rule $P = |\psi|^2$?
  \item Why complex Hilbert space (not real or quaternionic)?
  \item What is measurement, ontologically?
  \item What is entanglement?
  \item Why $[\hat{x}, \hat{p}] = i\hbar$?
\end{enumerate}

Hardy~\cite{Hardy2001} and Chiribella~\emph{et al.}~\cite{Chiribella2011}
have reconstructed QM from information-theoretic axioms, but such
reconstructions leave open: why \emph{these} axioms?  What forces them?

Recognition Science~\cite{RS_Axioms} provides a complete answer:
quantum mechanics is \emph{forced} by cost minimization from a single
primitive---the Recognition Composition Law.

%% ============================================================
\section{The Single Primitive}
\label{sec:primitive}
%% ============================================================

\begin{definition}[Recognition Composition Law]
The cost of composing recognition events satisfies a symmetric,
normalized defect measure.  For a cost function $F$ on positive ratios,
the RCL states:
\begin{equation}
  F(xy) + F(x/y) = 2F(x)F(y) + 2F(x) + 2F(y).
  \label{eq:RCL}
\end{equation}
\end{definition}

Together with three constraints:
\begin{enumerate}[label=(A\arabic*)]
  \item \textbf{Normalization}: $F(1) = 0$.
  \item \textbf{Symmetry}: $F(x) = F(1/x)$ for all $x > 0$.
  \item \textbf{Calibration}: $F''(1) = 1$.
\end{enumerate}

\begin{theorem}[Uniqueness of the Cost Functional]
\label{thm:unique-J}
The unique continuous function satisfying~\eqref{eq:RCL} and (A1)--(A3) is
\begin{equation}
  J(x) = \frac{1}{2}\left(x + x^{-1}\right) - 1.
\end{equation}
\end{theorem}

\begin{proof}[Proof sketch]
Setting $f(t) = F(e^t) + 1$ and substituting $x = e^s$, $y = e^t$ in
\eqref{eq:RCL} gives the d'Alembert equation
$f(s+t) + f(s-t) = 2f(s)f(t)$.  By the Acz\'el classification
theorem~\cite{Aczel1966}, the unique continuous solution with $f(0) = 1$
(from A1) and $f''(0) = 1$ (from A3) is $f(t) = \cosh t$, giving
$F(x) = \tfrac{1}{2}(x + x^{-1}) - 1$.  Symmetry (A2) is then
automatically satisfied.
\end{proof}

%% ============================================================
\section{The Forcing Chain T0--T8}
\label{sec:chain}
%% ============================================================

Each level is \emph{forced}, not postulated.

\subsection{T0: Logic from Cost}

\begin{remark}[Logic from Cost --- Domain Observation]
The domain of $J$ is $(0, \infty)$.  For $x \leq 0$, $x^{-1}$ is
undefined or negative, giving no valid cost assignment.  In this sense,
the domain of $J$ is identified with the space of consistent states.
This is a structural identification: consistent configurations are those
with well-defined $J$-cost, and contradictions (states with $x \leq 0$)
fall outside the domain.  The RS claim that this derives ``logic from
cost'' means: the logical constraint that states be consistent is
equivalent to the requirement that $J$-cost be finite.
\end{remark}

\subsection{T1: Existence Over Non-Existence}

\begin{proposition}[Meta-Principle --- Structural Argument]
In the RS identification of $x \to 0^+$ with non-existence:
$J(x) \to \infty$ (infinite cost for the zero-ratio limit).
Only $x = 1$ has zero cost.  Existence ($x = 1$) is structurally
preferred over non-existence ($x \to 0^+$) because its cost is minimal.
\end{proposition}

\begin{proof}
$J(x) = \tfrac{1}{2}(x + x^{-1}) - 1$.  As $x \to 0^+$,
$x^{-1} \to +\infty$, so $J(x) \to +\infty$.  At $x = 1$:
$J(1) = 0$.  Existence is cheaper than non-existence.
\end{proof}

\subsection{T2: Discreteness}

\begin{proposition}[Discreteness Forced --- Structural Argument]
Continuous configurations converge to the unique minimum $x = 1$ under
$J$-gradient flow.  In a system with self-similar composition, the
$\varphi$-ladder $\{\varphi^r : r \in \mathbb{Z}\}$ is the unique
discrete set closed under both $x \mapsto \varphi x$ and $x \mapsto x^{-1}$.
\end{proposition}

\begin{proof}[Structural argument]
Under $J$-gradient flow, any $x \neq 1$ is driven toward $x = 1$
(the unique global minimum).  The self-similar composition law
$\varphi^2 = \varphi + 1$ (from T6) identifies the natural discrete
scale as the $\varphi$-ladder.  Configurations off the ladder are
unstable under $J$-cost minimization.

\emph{Circular dependency note}: The identification of the
$\varphi$-ladder here relies on the golden ratio $\varphi$ from
T6, which formally appears later in the forcing chain.  The forcing
chain is therefore not strictly sequential at this step: T2 (discreteness)
uses the output of T6 ($\varphi$).  The correct interpretation is that
T2 and T6 are mutually reinforcing: self-similar composition forces
discreteness, and the unique ratio satisfying self-similar composition
is $\varphi$.
\end{proof}

\subsection{T3: Ledger (Double-Entry Bookkeeping)}

\begin{proposition}[Ledger Structure --- Structural Identification]
The symmetry $J(x) = J(1/x)$ motivates double-entry bookkeeping:
every recognition event $x$ has a matching inverse $x^{-1}$ with
equal cost, providing the RS origin of conservation laws.
\end{proposition}

\begin{proof}
$J(x) = \tfrac{1}{2}(x + x^{-1}) - 1 = \tfrac{1}{2}(x^{-1} + x) - 1
= J(x^{-1})$.  The cost cannot distinguish $x$ from $x^{-1}$;
both must be recorded.
\end{proof}

\subsection{T4: Recognition Operator}

\begin{proposition}[Recognition Operator --- Structural Identification]
A discrete operator $\widehat{R}$ is identified as the fundamental
evolution operator: it extracts observables from ledger states, enforces
$J$-minimization, and maintains ledger consistency.  $\widehat{R}$
replaces the Hamiltonian $\widehat{H}$.
\end{proposition}

\begin{proof}[Structural argument]
In the discrete ledger (T2), a physical quantity is observable only
when it corresponds to a stabilized recognition event---a ledger entry
with definite $J$-cost.  The map from ledger states to observables
must: (i)~satisfy $J$-minimization (the cost is non-increasing under
the dynamics), (ii)~preserve the ledger double-entry structure (T3),
and (iii)~be realizable in the $2^D = 8$-tick period (T7).  These
three conditions uniquely specify a linear discrete-time operator on
the Hilbert space of ledger states, which we call $\widehat{R}$.
$\widehat{R}$ is a discrete analog of the time-evolution operator
$e^{-i\widehat{H}\tau_0/\hbar}$, but derived from the $J$-cost structure
rather than postulated.
\end{proof}

Dynamics are given by the discrete update:
\begin{equation}
  s(t + 8\tau_0) = \widehat{R}\bigl(s(t)\bigr),
\end{equation}
minimizing $J$-cost rather than energy.

\subsection{T5: Unique $J$}

The cost functional $J(x) = \tfrac{1}{2}(x + x^{-1}) - 1$ is uniquely
determined by the RCL plus normalization and calibration
(Theorem~\ref{thm:unique-J}).

\subsection{T6: The Golden Ratio $\varphi$}

\begin{theorem}[$\varphi$ Forced]
Self-similar nesting in the discrete ledger requires $x^2 = x + 1$.
The unique positive solution is $\varphi = (1+\sqrt{5})/2$.
\end{theorem}

\begin{proof}
The self-similar condition requires that the ratio between successive
levels of the hierarchy satisfies $x = 1 + 1/x$, i.e., $x^2 = x + 1$.
Solving: $x = (1 \pm \sqrt{5})/2$.  The positive root is $\varphi$.
\end{proof}

\subsection{T7: The 8-Tick Epoch}

\begin{proposition}[8-Tick Epoch --- Structural Argument]
The linking structure of the discrete ledger is identified as requiring
a Hamiltonian cycle on the $D$-cube graph.  The minimal Hamiltonian
cycle length on $Q_D$ (which has $2^D$ vertices) is $2^D$.
\end{proposition}

\begin{proof}[Structural argument]
The discrete ledger in $D$ spatial dimensions has $2^D$ states
(one per sign pattern of the $D$ spatial directions).  The
recognition cycle must visit all $2^D$ states to ensure complete
ledger updating---a Hamiltonian cycle on the $D$-cube graph
$Q_D$.  The $D$-cube graph is bipartite (vertices alternately
colored black and white by parity); the shortest Hamiltonian cycle
on $Q_D$ has length $2^D$ (it visits all vertices once and returns).
For $D = 3$: $2^3 = 8$ ticks.
\end{proof}

\begin{remark}
T7 and T8 are mutually reinforcing: the ledger-linking structure forces
$n = 8$ as the minimal period (from the unique non-trivial linking in
$D = 3$), while $n = 8$ forces $D = 3$ via $2^D = 8$ (uniquely).
The correct causal direction is: the linking structure of the
recognition ledger forces the minimal period $n = 8$, and this forces
$D = 3$.  The 8-tick epoch defines the fundamental time unit $\tau_0$
and the coherence energy $E_{\mathrm{coh}} = \varphi^{-5}$.
\end{remark}

\subsection{T8: Three Spatial Dimensions}

\begin{proposition}[$D = 3$ Forced]
The equation $2^D = 8$ (from T7) has the unique positive-integer
solution $D = 3$.  This is confirmed by two independent structural
requirements: (i)~nontrivial topological linking (knots exist only in
$D = 3$), and (ii)~gap-45 synchronization
($\mathrm{lcm}(8, 45) = 360$, yielding angular closure).
\end{proposition}

\begin{proof}
Primary: $2^D = 8$ has the unique positive-integer solution $D = 3$
(since $2^1 = 2$, $2^2 = 4$, $2^3 = 8$, and $2^D$ is strictly
increasing).
Confirmation~(i): nontrivial knots require $D = 3$ (in $D \neq 3$,
curves are either trivially linked or cannot be knotted).  The
recognition ledger requires nontrivial linking to maintain closure
of the 8-tick walk; hence $D \geq 3$.  Combined with $2^D = 8$,
this forces $D = 3$.
Confirmation~(ii): the gauge angle structure has a 45-unit gap (the
Gap-45 synchronization); for angular closure, the total angular
period must be divisible by both $8\tau_0$ (the tick period) and
$45$ (the gap); $\mathrm{lcm}(8, 45) = 360$, which is the full
circle in degrees---confirming $D = 3$ independently.
\end{proof}

%% ============================================================
\section{Quantum Mechanics Emerges}
\label{sec:QM}
%% ============================================================

\subsection{Complex Hilbert Space}

\begin{proposition}[Complex Numbers --- Structural Argument]
The 8-tick epoch has half-period 4 ticks, corresponding to a phase
rotation of $\pi$.  The RS identification is that representing this
requires a number system in which $e^{i\pi} = -1$ has a solution:
$\mathbb{C}$.
\end{proposition}

\begin{proof}
$\mathbb{R}$ cannot represent $e^{i\pi} = -1$ internally (there is
no element $i$ with $i^2 = -1$).  $\mathbb{H}$ (quaternions) would require the phase structure to accommodate
all 8 generators $\{\pm 1, \pm i, \pm j, \pm k\}$ of the quaternion group
$Q_8$, which has 3 independent imaginary units --- more structure than
the single phase rotation $e^{i\pi}$ of the 8-tick half-period.  In the
RS framework, the 8-tick complex phase $e^{2\pi i k/8}$ for $k = 0,\ldots,7$
generates $\mathbb{Z}_8 \subset \mathbb{C}^\times$; this is naturally a
subgroup of $\mathbb{C}^\times$, not of the quaternion group.  Only
$\mathbb{C}$ provides the minimal number system accommodating both the
half-period $e^{i\pi} = -1$ and the cyclic 8-tick structure.
(This argument is a structural identification rather than a proof of
uniqueness; a complete derivation from the RS Hamiltonian is an
identified open problem.)
\end{proof}

\begin{corollary}[Euler's Identity]
$e^{i\pi} + 1 = 0$ is a structural necessity: $i$ is the square root
of the half-epoch operator, and Euler's identity records that the
half-epoch returns the state to its negative.
\end{corollary}

\subsection{The Born Rule}

\begin{proposition}[Born Rule from Phase Invariance and Gleason's Theorem]
The probability of a measurement outcome is $P = |A|^2$, where $A$ is
the amplitude.
\end{proposition}

\begin{proof}
The RS recognition rate $R(Ae^{i\theta})$ for a state with amplitude
$A$ and phase $\theta$ must satisfy:
\begin{enumerate}[label=(\roman*)]
  \item Phase invariance: $R(Ae^{i\theta}) = R(A)$ (the recognition
  rate cannot depend on the gauge-choice phase).
  \item Subadditivity: $R(A+B) \leq R(A) + R(B)$.
  \item Normalization: $R(1) = 1$.
\end{enumerate}
By~(i), $R$ depends only on $|A|$, so $R(A) = f(|A|)$ for some
function $f: \mathbb{R}_{\geq 0} \to \mathbb{R}_{\geq 0}$.
By~(iii), $f(1) = 1$.  By Gleason's theorem~\cite{Gleason1957}
(which requires $\dim \geq 3$, ensured by T8), any frame function
on a Hilbert space of dimension $\geq 3$ must take the form
$f(r) = r^2$, i.e., $R(A) = |A|^2$.

\emph{Note on the subadditivity condition}: The condition (ii) as stated
(subadditivity) is weaker than Gleason's actual hypotheses, which require
the probability assignment to be \emph{additive} on orthogonal subspaces.
The RS structural motivation for additivity is that orthogonal recognition
channels are independent and their combined $J$-cost is additive;
subadditivity follows from additivity.  The full derivation of the
additivity condition from the RS double-entry structure is an identified
open problem; here we appeal to Gleason's theorem with the RS confirmation
that $D = 3$ ensures $\dim \geq 3$.
\end{proof}

\subsection{Entanglement}

\begin{definition}[Entanglement as Shared Constraint]
Two recognition events $A$ and $B$ are entangled when their joint
ledger state has a conserved quantity $Q_{AB} = Q_A + Q_B = q_0 =$
const.  This forces correlations: knowing $Q_A$ instantly determines
$Q_B = q_0 - Q_A$.
\end{definition}

\begin{remark}
No action at a distance is required.  The correlation is in the ledger
entry, not in a signal traveling between $A$ and $B$.  Bell inequality
violation occurs because the shared ledger constraint creates a
non-factorizable $J$-cost that no local hidden variable model can
reproduce.
\end{remark}

\subsection{Commutation Relations}

\begin{proposition}[Canonical Commutation --- Structural Argument]
$[\hat{x}, \hat{p}] = i\hbar$, where $\hbar = E_{\mathrm{coh}}\tau_0
= \varphi^{-5}\tau_0$.
\end{proposition}

\begin{proof}[Proof sketch]
The ledger's double-entry structure (T3) makes the order of
observation matter: a position measurement at tick $t$ commits a
ledger entry (setting $x$ to a definite value), which back-reacts on
the momentum state through the recognition update at $t+1$.  The
canonical commutation relation arises from the discrete group structure
of the ledger updates: the translation operators $\hat{T}_x(\Delta x)
= e^{i\Delta x\,\hat{p}/\hbar}$ and $\hat{T}_p(\Delta p)
= e^{-i\Delta p\,\hat{x}/\hbar}$ satisfy
\[
  \hat{T}_x(\Delta x)\hat{T}_p(\Delta p)
  = e^{i\Delta x\,\Delta p/\hbar}\,
    \hat{T}_p(\Delta p)\hat{T}_x(\Delta x),
\]
giving $[\hat{x}, \hat{p}] = i\hbar$ in the infinitesimal limit.
The coherence quantum $\hbar = E_{\rm coh}\tau_0 = \varphi^{-5}\tau_0$
sets the scale.  A complete derivation showing that the RS recognition
operator $\widehat{R}$ generates exactly this Weyl--Heisenberg algebra
is available in the detailed RS architecture; the above sketch
establishes the structural form of the commutator.
\end{proof}

\subsection{The Measurement Problem: Dissolved}

\begin{proposition}[Measurement is Recognition --- Structural Identification]
In RS, measurement is identified with a recognition event.  The ``collapse''
of the wavefunction is the discrete update of the ledger at the
completion of the 8-tick epoch.  Under this identification, no
additional measurement postulate is needed.
\end{proposition}

\begin{proof}
The RS ontology distinguishes the ledger (physical reality) from
$\psi$ (a computational description of $J$-cost-weighted probability
distributions over ledger states).  Superposition exists in $\psi$,
not in the ledger.  A recognition event (measurement) produces a new
ledger entry, which is either present (event occurred) or absent
(did not occur)---no superposition in the ledger.  The ``collapse'' is
the updating of $\psi$ to match the observed ledger entry.
\end{proof}

\subsection{Decoherence}

The decoherence time for a system with energy uncertainty $\Delta E$ is
\begin{equation}
  \tau_D = \frac{\tau_0}{\Delta E/E_{\mathrm{coh}}}
  = \frac{E_{\mathrm{coh}}\,\tau_0}{\Delta E}
  = \frac{\hbar}{\Delta E},
\end{equation}
since $\hbar = E_{\mathrm{coh}}\,\tau_0 = \varphi^{-5}\,\tau_0$.
For macroscopic objects ($\Delta E \gg E_{\mathrm{coh}}$),
$\tau_D \ll \tau_0$---decoherence is effectively instantaneous,
consistent with classical behavior.

\subsection{Interpretation Resolution}

\begin{center}
\renewcommand{\arraystretch}{1.3}
\begin{tabular}{@{}lp{8cm}@{}}
\toprule
\textbf{Interpretation} & \textbf{RS Assessment} \\
\midrule
Copenhagen & Correct operationally: ``don't ask what $\psi$ is'' works
because $\psi$ is a tool, not ontology. \\
Many-Worlds & Incorrect: there is one ledger, not many branches. \\
Bohmian & Partially correct: the pilot wave is the $J$-cost gradient
field. \\
Relational & Correct at the price of circularity: RS provides the
non-relational foundation (the ledger). \\
\bottomrule
\end{tabular}
\end{center}

%% ============================================================
\section{The Quantum-to-Classical Transition}
\label{sec:classical}
%% ============================================================

As the number of ledger cells $N \to \infty$, the law of large numbers
gives $\langle x \rangle_N \to E[x]$.  The key observation is that for
macroscopic systems, fluctuations of $x_i$ around their mean become small
(relative fluctuation $\sim N^{-1/2}$), so the $J$-cost average is well
approximated by $J$ evaluated at the mean:
\begin{equation}
  \langle J \rangle_N = E[J(x)] \approx J(E[x]) + O(N^{-1}).
\end{equation}
(By Jensen's inequality, $J(E[x]) \leq E[J(x)]$ with equality only
when all $x_i$ are equal---i.e., in the classical deterministic limit.)
In the large-$N$ limit, individual cells track the mean, and dynamics
become deterministic.  This is the classical limit: quantum fluctuations
of order $\hbar/(N\,\text{action})$ are negligible for macroscopic~$N$.

%% ============================================================
\section{Conclusion}
\label{sec:conclusion}
%% ============================================================

Quantum mechanics is derived from cost minimization.  The single
primitive---the Recognition Composition Law---forces the unique cost
function $J$, the 8-tick discrete dynamics, and the recognition
operator $\widehat{R}$.  The results:
\begin{enumerate}
  \item Complex Hilbert space: forced by the 8-tick half-period
  $e^{i\pi} = -1$.
  \item Born rule $P = |\psi|^2$: from phase invariance,
  subadditivity, and normalization.
  \item Entanglement: from shared ledger constraints.
  \item $[\hat{x}, \hat{p}] = i\hbar$: from double-entry
  non-commutativity, with $\hbar = \varphi^{-5}\tau_0$.
  \item Measurement collapse: from the discrete ledger update at
  $C \geq 1$.
\end{enumerate}

No measurement postulate, no Hilbert-space axiom, and no free
parameters.  Quantum mechanics is \emph{forced}, not postulated.

%% ============================================================
\begin{thebibliography}{99}

\bibitem{Aczel1966}
J.~Acz\'el, \emph{Lectures on Functional Equations and Their Applications},
Academic Press, New York (1966).

\bibitem{Hardy2001}
L.~Hardy, ``Quantum Theory from Five Reasonable Axioms,'' arXiv:
quant-ph/0101012 (2001).

\bibitem{Chiribella2011}
G.~Chiribella, G.~M.~D'Ariano, and P.~Perinotti, ``Informational
Derivation of Quantum Theory,'' Phys.\ Rev.\ A \textbf{84}, 012311
(2011).

\bibitem{Gleason1957}
A.~M.~Gleason, ``Measures on the Closed Subspaces of a Hilbert
Space,'' J.\ Math.\ Mech.\ \textbf{6}, 885 (1957).

\bibitem{Zurek2003}
W.~H.~Zurek, ``Decoherence, Einselection, and the Quantum Origins
of the Classical,'' Rev.\ Mod.\ Phys.\ \textbf{75}, 715 (2003).

\bibitem{RS_Axioms}
J.~Washburn, ``The Algebra of Reality: A Recognition Science Derivation
of Physical Law,'' \emph{Axioms} \textbf{15}(2), 90 (2026).
\texttt{doi:10.3390/axioms15020090}

\end{thebibliography}

\end{document}
